\documentclass[11pt]{memoir}
\setcounter{secnumdepth}{4}
\setcounter{tocdepth}{3}

% ==========================================
% FONTS: EB Garamond (matching Volume I)
% ==========================================
\usepackage[T1]{fontenc}
\usepackage[utf8]{inputenc}
\usepackage{lmodern}
\frenchspacing

\usepackage[
  cmintegrals,
  cmbraces,
  noamssymbols
]{newtxmath}
\usepackage{ebgaramond}

\usepackage[
  activate={true,nocompatibility},
  final,
  tracking=true,
  kerning=true,
  spacing=true,
  factor=1100,
  stretch=10,
  shrink=10
]{microtype}

% ==========================================
% PACKAGES
% ==========================================
\usepackage{mleftright}
\mleftright
\let\Bbbk\undefined
\let\openbox\undefined
\usepackage{amsmath,amssymb,amsthm}
\usepackage{mathtools}
\mathtoolsset{showonlyrefs,showmanualtags}
\usepackage{tikz-cd}
\usetikzlibrary{positioning,arrows.meta,calc}
\makeatletter
\def\hyper@nopatch@caption{}
\makeatother
\usepackage[unicode,pdfencoding=auto,psdextra]{hyperref}
\usepackage{geometry}
\geometry{top=1.2in, bottom=1.2in, left=1.25in, right=1.25in, footskip=0.5in}
\setlength{\headheight}{28pt}
\setlength{\emergencystretch}{3em}
\setlength{\hfuzz}{1pt}
\hbadness=3000
\raggedbottom
\usepackage{enumitem}
\usepackage{longtable}
\usepackage{multicol}
\usepackage{tcolorbox}

% ==========================================
% CROSS-REFERENCES TO VOLUME I
% ==========================================
% The xr-hyper package allows referencing labels from the main volume.
% Build Vol I first so that main.aux exists; then build the concordance.
\usepackage{xr-hyper}
\externaldocument{../main}

% ==========================================
% THEOREM ENVIRONMENTS (matching Vol I)
% ==========================================
\usepackage{thmtools}

\declaretheoremstyle[
  spaceabove=4pt plus 1pt minus 1pt,
  spacebelow=4pt plus 1pt minus 1pt,
  headfont=\normalfont\scshape,
  notefont=\normalfont\itshape,
  bodyfont=\normalfont,
  postheadspace=0.5em,
  headpunct={.}
]{garamondthm}

\declaretheoremstyle[
  spaceabove=4pt plus 1pt minus 1pt,
  spacebelow=4pt plus 1pt minus 1pt,
  headfont=\normalfont\itshape,
  notefont=\normalfont\itshape,
  bodyfont=\normalfont,
  postheadspace=0.5em,
  headpunct={.}
]{garamonddef}

\newcommand{\ClaimStatusProvedHere}{\textnormal{[proved here]}}
\newcommand{\ClaimStatusProvedElsewhere}{\textnormal{[proved elsewhere]}}
\newcommand{\ClaimStatusOpen}{\textnormal{[open]}}
\newcommand{\ClaimStatusConjectured}{\textnormal{[conjectured]}}
\newcommand{\ClaimStatusHeuristic}{\textnormal{[physical heuristic]}}

\declaretheorem[style=garamondthm, name=Theorem, numberwithin=section]{theorem}
\declaretheorem[style=garamondthm, name=Lemma, sibling=theorem]{lemma}
\declaretheorem[style=garamondthm, name=Proposition, sibling=theorem]{proposition}
\declaretheorem[style=garamondthm, name=Corollary, sibling=theorem]{corollary}
\declaretheorem[style=garamondthm, name=Verification, sibling=theorem]{verification}
\declaretheorem[style=garamondthm, name=Computation, sibling=theorem]{computation}
\declaretheorem[style=garamondthm, name=Main Theorem, sibling=theorem]{maintheorem}
\declaretheorem[style=garamonddef, name=Definition, sibling=theorem]{definition}
\declaretheorem[style=garamonddef, name=Example, sibling=theorem]{example}
\declaretheorem[style=garamonddef, name=Remark, sibling=theorem]{remark}
\declaretheorem[style=garamonddef, name=Conjecture, sibling=theorem]{conjecture}
\declaretheorem[style=garamonddef, name=Notation, sibling=theorem]{notation}
\declaretheorem[style=garamonddef, name=Convention, sibling=theorem]{convention}
\declaretheorem[style=garamonddef, name=Construction, sibling=theorem]{construction}
\declaretheorem[style=garamonddef, name=Spectral Sequence, sibling=theorem]{spectralsequence}
\declaretheorem[style=garamonddef, name=Calculation, sibling=theorem]{calculation}
\declaretheorem[style=garamonddef, name=Technique, sibling=theorem]{technique}
\declaretheorem[style=garamonddef, name=Perspective, sibling=theorem]{perspective}
\declaretheorem[style=garamonddef, name=Summary, sibling=theorem]{summary}
\declaretheorem[style=garamonddef, name=Observation, sibling=theorem]{observation}
\declaretheorem[style=garamonddef, name=Open Problem, sibling=theorem]{openproblem}
\declaretheorem[style=garamonddef, name=Strategy, sibling=theorem]{strategy}
\declaretheorem[style=garamonddef, name=Framework, sibling=theorem]{framework}
\declaretheorem[style=garamonddef, name=Context, sibling=theorem]{context}
\declaretheorem[style=garamonddef, name=Historical, sibling=theorem]{historical}

\makeatletter
\def\theHtheorem{\thesection.\arabic{theorem}}
\let\theHlemma\theHtheorem
\let\theHproposition\theHtheorem
\let\theHcorollary\theHtheorem
\let\theHverification\theHtheorem
\let\theHcomputation\theHtheorem
\let\theHdefinition\theHtheorem
\let\theHexample\theHtheorem
\let\theHremark\theHtheorem
\let\theHconjecture\theHtheorem
\let\theHnotation\theHtheorem
\let\theHconvention\theHtheorem
\let\theHconstruction\theHtheorem
\let\theHmaintheorem\theHtheorem
\let\theHperspective\theHtheorem
\let\theHsummary\theHtheorem
\let\theHobservation\theHtheorem
\let\theHopenproblem\theHtheorem
\let\theHstrategy\theHtheorem
\let\theHframework\theHtheorem
\let\theHcontext\theHtheorem
\let\theHhistorical\theHtheorem
\let\theHspectralsequence\theHtheorem
\let\theHcalculation\theHtheorem
\let\theHtechnique\theHtheorem
\def\theHtable{\theHchapter.\arabic{table}}
\def\theHfigure{\theHchapter.\arabic{figure}}
\makeatother

\newenvironment{evidence}[1][Evidence]{%
  \par\medskip\noindent\textit{#1.}\enspace\ignorespaces
}{%
  \par\medskip
}

% ==========================================
% MATH OPERATORS AND NOTATION (matching Vol I)
% ==========================================
\DeclareMathOperator{\Hom}{Hom}
\DeclareMathOperator{\RHom}{RHom}
\DeclareMathOperator{\End}{End}
\DeclareMathOperator{\Ext}{Ext}
\DeclareMathOperator{\Tor}{Tor}
\DeclareMathOperator{\Res}{Res}
\DeclareMathOperator{\Ind}{Ind}
\DeclareMathOperator{\colim}{colim}
\DeclareMathOperator{\Ran}{Ran}
\DeclareMathOperator{\sgn}{sgn}
\DeclareMathOperator{\Free}{Free}
\DeclareMathOperator{\Cofree}{Cofree}
\DeclareMathOperator{\Spec}{Spec}

\newcommand{\Com}{\mathsf{Com}}
\newcommand{\Lie}{\mathsf{Lie}}
\newcommand{\Ass}{\mathsf{Ass}}
\newcommand{\Pois}{\mathsf{Pois}}
\newcommand{\BV}{\mathsf{BV}}
\newcommand{\Grav}{\mathsf{Grav}}
\newcommand{\Eone}{\mathsf{E}_1}
\newcommand{\Etwo}{\mathsf{E}_2}
\newcommand{\En}{\mathsf{E}_n}
\newcommand{\Einf}{\mathsf{E}_{\infty}}
\newcommand{\Linf}{\mathsf{L}_{\infty}}
\newcommand{\Ainf}{\mathsf{A}_{\infty}}
\newcommand{\Cinf}{\mathsf{C}_{\infty}}
\newcommand{\Pinf}{\mathsf{P}_{\infty}}

\newcommand{\chirLie}{\mathsf{Lie}^{\mathrm{ch}}}
\newcommand{\chirAss}{\mathsf{Ass}^{\mathrm{ch}}}
\newcommand{\chirCom}{\mathsf{Com}^{\mathrm{ch}}}
\newcommand{\chirPois}{\mathsf{Pois}^{\mathrm{ch}}}

\newcommand{\Chir}{\mathsf{Chir}}
\newcommand{\Fact}{\mathsf{Fact}}
\newcommand{\FM}{\mathrm{FM}}

\newcommand{\cat}[1]{\mathsf{#1}}
\newcommand{\Vect}{\cat{Vect}}
\newcommand{\Ch}{\cat{Ch}}
\newcommand{\Alg}{\cat{Alg}}
\newcommand{\CoAlg}{\cat{CoAlg}}
\newcommand{\Op}{\cat{Op}}
\newcommand{\DGCat}{\cat{DGCat}}
\newcommand{\DMod}{\cat{D}\text{-}\cat{Mod}}
\newcommand{\DX}{\mathcal{D}_X}
\newcommand{\OX}{\mathcal{O}_X}

\newcommand{\C}{\mathbb{C}}
\newcommand{\Z}{\mathbb{Z}}
\newcommand{\R}{\mathbb{R}}

\newcommand{\cA}{\mathcal{A}}
\newcommand{\cB}{\mathcal{B}}
\newcommand{\cV}{\mathcal{V}}
\newcommand{\barB}{\overline{B}}
\newcommand{\barPi}{\overline{\Pi}}

\newcommand{\dzero}{d_0}
\newcommand{\dfib}{d_{\mathrm{fib}}}
\newcommand{\Dg}[1]{D_{#1}}
\newcommand{\Dtot}{D_{\mathrm{tot}}}
\newcommand{\mcurv}[1]{m_0^{(#1)}}
\newcommand{\Defcyc}{\operatorname{Def}_{\mathrm{cyc}}}
\newcommand{\Gmod}{\mathcal{G}^{\mathrm{mod}}}
\newcommand{\Perbr}{\operatorname{Per}}
\newcommand{\ChirHoch}{CH}
\newcommand{\Walg}{\mathcal{W}}
\newcommand{\critLevel}{-h^\vee}

\newcommand{\BDref}[1]{\textup{[\cite{BD04}, #1]}}

\providecommand{\Aut}{\operatorname{Aut}}
\providecommand{\B}{\bar{B}}
\providecommand{\Cobar}{\Omega}
\providecommand{\gr}{\operatorname{gr}}
\providecommand{\id}{\mathrm{id}}
\providecommand{\ad}{\operatorname{ad}}
\providecommand{\dlog}{d\log}
\providecommand{\normord}[1]{:{#1}:}
\providecommand{\bA}{\mathbb{A}}
\providecommand{\bC}{\mathbb{C}}
\providecommand{\bR}{\mathbb{R}}
\providecommand{\bZ}{\mathbb{Z}}
\providecommand{\bQ}{\mathbb{Q}}
\providecommand{\cC}{\mathcal{C}}
\providecommand{\cD}{\mathcal{D}}
\providecommand{\cH}{\mathcal{H}}
\providecommand{\cM}{\mathcal{M}}
\providecommand{\cO}{\mathcal{O}}
\providecommand{\cP}{\mathcal{P}}
\providecommand{\cN}{\mathcal{N}}
\providecommand{\cL}{\mathcal{L}}
\providecommand{\fg}{\mathfrak{g}}
\providecommand{\VD}{\mathbb{D}}
\providecommand{\Map}{\operatorname{Map}}
\providecommand{\holim}{\operatorname{holim}}
\providecommand{\hocolim}{\operatorname{hocolim}}
\providecommand{\fib}{\operatorname{fib}}
\providecommand{\cofib}{\operatorname{cofib}}
\providecommand{\CE}{\mathrm{CE}}

\newcommand{\Ydg}{\mathcal{Y}^{\mathrm{dg}}}
\newcommand{\Factord}{\Fact_{\Eone}^{\mathrm{ord}}}

\usepackage{setspace}
\setsecheadstyle{\normalfont\Large\scshape}
\setsubsecheadstyle{\normalfont\large\scshape}

% ==========================================
% DOCUMENT
% ==========================================

\title{\textit{Concordance with Primary Literature}\\[0.5em]
\large Companion to \textit{Modular Homotopy Theory for Factorization Algebras on Curves}}
\author{Raeez Lorgat}
\date{March 2026}

\begin{document}

\maketitle

\begin{abstract}
\itshape\noindent
This document is the normative status ledger for the monograph
\textit{Modular Homotopy Theory for Factorization Algebras on Curves}
(Volumes~I and~II).  It records the precise relationship of the
monograph's constructions and results to the primary literature
(Beilinson--Drinfeld, Francis--Gaitsgory, Gui--Li--Zeng,
Loday--Vallette), defines the resolved PBW entry theorem and the
universal $\Theta_{\cA}$ package, stratifies the conjectural frontier
(MC3/MC4/MC5), and articulates the programme toward modular homotopy
theory for factorization algebras on curves.

\medskip\noindent
Cross-references to Volume~I are resolved via the \texttt{xr-hyper}
package.  Build Volume~I first to generate \texttt{main.aux}; then
build this document.
\end{abstract}

\tableofcontents

\chapter{Concordance with primary literature}
\label{chap:concordance}
\index{concordance|textbf}

This chapter establishes precise relationships between the constructions
of this monograph and the foundational references.  We begin with a summary
of the monograph's principal contributions (\S\ref{sec:principal-contributions}),
then provide detailed comparisons with each primary reference.


\section{Principal contributions}\label{sec:principal-contributions}

We isolate the aspects of this work that go beyond existing literature.

\begin{enumerate}[label=\emph{(\arabic*)}]

\item \emph{Configuration space bar-cobar adjunction
  (Theorems~\ref{thm:bar-cobar-isomorphism-main}
  and~\ref{thm:higher-genus-inversion}).}
  The bar and cobar constructions for chiral algebras are realized
  via explicit integrals over Fulton--MacPherson compactifications
  $\overline{C}_n(X)$, with the bar differential given by residues
  at collision divisors and the cobar differential by inclusions
  of strata.  The adjunction is proved for all Koszul chiral algebras
  and all genera.  This is the first complete proof of chiral Koszul
  duality via configuration space geometry.

\item \emph{Genus universality
  (Theorem~\ref{thm:genus-universality}).}
  For any Koszul chiral algebra~$\cA$, the genus-$g$ obstruction
  factors as $\mathrm{obs}_g(\cA) = \kappa(\cA) \cdot \lambda_g$,
  where $\kappa$ is an algebra-dependent constant and $\lambda_g$
  is the universal Faber--Pandharipande tautological class
  on~$\overline{\mathcal{M}}_g$.
  The genus-independent coefficient $\kappa(\cA)$ is determined
  by the genus-$1$ curvature.  This result is new.

\item \emph{Deformation-obstruction complementarity
  (Theorem~\ref{thm:quantum-complementarity-main}).}
  For a Koszul pair $(\cA, \cA^!)$, the genus-$g$ obstructions
  satisfy $\mathrm{obs}_g(\cA) + \mathrm{obs}_g(\cA^!) =
  [\text{characteristic class of } Z(\cA)]$ in $H^*(\overline{\mathcal{M}}_g)$.
  The Kodaira--Spencer map is constructed for all Koszul pairs
  and all genera (Theorem~\ref{thm:kodaira-spencer-chiral-complete}).
  This is the first formulation and proof of complementarity
  for chiral algebras.

\item \emph{$\Eone$-chiral Koszul duality
  (Theorem~\ref{thm:e1-chiral-koszul-duality}).}
  The bar-cobar adjunction extends to $\Eone$-chiral algebras
  (nonlocal vertex algebras with $R$-matrix braiding), covering
  Yangians and toroidal algebras.  This goes beyond the
  $\Einf$ (commutative/local) setting of BD and FG.

\item \emph{Explicit computations.}
  Complete Koszul duality computations for all standard families:
  Heisenberg, free fermion, $bc$-$\beta\gamma$, affine Kac--Moody
  ($\widehat{\fg}_k \leftrightarrow \widehat{\fg}_{-k-2h^\vee}$),
  Virasoro ($\mathrm{Vir}_c \leftrightarrow \mathrm{Vir}_{26-c}$),
  $\mathcal{W}_N$ algebras, Yangians.
  Genus-$g$ free energies computed through $g = 5$ with
  Bernoulli asymptotics.  These explicit results are new.

\end{enumerate}

\section{Relationship to Beilinson--Drinfeld}

\begin{tabular}{>{\raggedright\arraybackslash}p{0.38\textwidth}>{\raggedright\arraybackslash}p{0.52\textwidth}}
\textbf{Our Terminology} & \textbf{BD Terminology} \\
\hline
Chiral algebra & Chiral algebra (same) \\
$\Einf$-chiral algebra & Vertex algebra / Chiral algebra \\
$\Eone$-chiral algebra & No direct analog (generalization) \\
Chiral bracket $\mu$ & Chiral operation $\mu: j_* j^*(A \boxtimes A) \to \Delta_! A$ \\
Bar complex $\B(\cA)$ & Chiral chain complex (cf.\ Chevalley--Eilenberg for Lie case) \\
Factorization structure & Factorization algebra structure \\
Ran space $\Ran(X)$ & $\mathfrak{R}(X)$ in BD notation \\
\end{tabular}

\begin{remark}[Key extension]
Our $\Eone$-chiral algebras extend BD's chiral algebras by dropping skew-symmetry. BD's chiral algebras are our $\Einf$-chiral algebras.
\end{remark}

\section{Relationship to Francis--Gaitsgory}
\index{Francis--Gaitsgory}

\begin{tabular}{>{\raggedright\arraybackslash}p{0.38\textwidth}>{\raggedright\arraybackslash}p{0.52\textwidth}}
\textbf{Our Terminology} & \textbf{FG Terminology} \\
\hline
Chiral Koszul duality & Chiral Koszul duality (same) \\
$\chirCom$-$\chirLie$ duality & Main theorem of FG \\
$\chirAss$-$\chirAss$ duality & Not explicitly treated (implicit) \\
Pro-nilpotence & Pro-nilpotent tensor $\infty$-category \\
Bar-cobar equivalence & Koszul duality equivalence \\
\end{tabular}

\begin{remark}[Key extension]
FG establish $\chirCom$-$\chirLie$ duality. We show this is derived from the more fundamental $\chirAss$-$\chirAss$ self-duality via the deformation Pois → Ass.
\end{remark}

\section{Relationship to Gui--Li--Zeng}
\index{Gui--Li--Zeng}

\begin{tabular}{>{\raggedright\arraybackslash}p{0.38\textwidth}>{\raggedright\arraybackslash}p{0.52\textwidth}}
\textbf{Our Terminology} & \textbf{GLZ Terminology} \\
\hline
Quadratic chiral algebra & Quadratic chiral algebra (same) \\
Koszul dual $\cA^!$ & Quadratic dual $A^!$ \\
Bar complex & Not used (direct quadratic dual) \\
Non-quadratic duality & Not treated \\
$\Eone$-chiral algebras & Not treated \\
\end{tabular}

\begin{proposition}[GLZ as special case; \ClaimStatusProvedHere]\label{prop:glz-special-case}
The Gui--Li--Zeng quadratic duality is a special case of the $\chirAss$-duality of this monograph:
\[
\text{GLZ duality} = \chirAss\text{-duality} |_{\Einf\text{-chiral}} |_{\text{quadratic}}
\]
Restricting to $\Einf$-chiral algebras with quadratic presentations recovers their theory.
\end{proposition}

\begin{proof}
An $\Einf$-chiral algebra $\cA$ with quadratic presentation has a quadratic relation ideal $R \subset j_*j^*(\cA^{\boxtimes 2})$.  Its $\chirAss$ Koszul dual is determined by bar degrees $\leq 2$: the generators $V^*$ in $\bar{B}^1$ and the orthogonal complement of relations $R^\perp$ in $\bar{B}^2$ (the higher bar degrees, while nonzero, contribute no independent data for quadratic algebras).  The Koszul dual $\cA^! = \Omega(\bar{B}(\cA))$ is then determined by $R^\perp$ under the chiral pairing $j_*j^*(\cA^{\boxtimes 2}) \otimes j_*j^*((\cA^!)^{\boxtimes 2}) \to \omega_X$.  This is precisely the quadratic dual construction of GLZ.

For the $\Einf$ restriction: our $\chirAss$ bar complex uses the full associative chiral product $\mu\colon j_*j^*(\cA \boxtimes \cA) \to \Delta_!\cA$.  When $\cA$ is $\Einf$-chiral (i.e., the chiral bracket is skew-symmetric), the bar complex acquires a commutative coalgebra structure, and the Koszul dual is a Lie chiral algebra.  This recovers the FG framework as a special case (see Theorem~\ref{thm:fg-from-assch}).
\end{proof}

\begin{theorem}[FG duality from $\chirAss$ self-duality; \ClaimStatusProvedHere]\label{thm:fg-from-assch}
The Francis--Gaitsgory $\chirCom$-$\chirLie$ duality is the associated graded of the $\chirAss$ self-duality under the PBW filtration.  Precisely: for a quadratic $\Einf$-chiral algebra $\cA$, the $\chirAss$ bar complex $\bar{B}_{\chirAss}(\cA)$ carries a filtration $F_\bullet$ (induced by the symmetrization degree of the $\Sigma$-action on $C_n(X)$) such that:
\begin{equation}\label{eq:fg-from-assch}
\mathrm{gr}^F \bar{B}_{\chirAss}(\cA) \;\cong\; \bar{B}_{\chirCom}(\mathrm{Sym}^{\mathrm{ch}}(\cA)) \;\cong\; \mathrm{CE}^{\chirLie}(\cA)
\end{equation}
and the associated spectral sequence degenerates at $E_1$ for Koszul algebras.
\end{theorem}

\begin{proof}
\emph{Step~1: PBW filtration.}
The chiral bar complex $\bar{B}^n_{\chirAss}(\cA)$ consists of sections on $\overline{C}_{n+1}(X)$ with coefficients in $\cA^{\boxtimes(n+1)} \otimes \Omega^n_{\log}$.  Define $F_p \bar{B}^n$ as the span of elements symmetric under at least $(n-p)$ transpositions of configuration space points.  This is the chiral analogue of the arity filtration on the classical associative bar complex (Loday--Vallette \cite{LV12}, \S9.3.2).

\emph{Step~2: Associated graded identification.}
On $\overline{C}_n(X)$, the $\Sigma_n$-representation decomposes as $j_*j^*(\cA^{\boxtimes n}) = \mathrm{Sym}^n_{\mathrm{ch}}(\cA) \oplus \bigwedge^2_{\mathrm{ch}}(\cA) \otimes \cdots$.  The PBW filtration separates these: $F_0/F_1 = \mathrm{Sym}^n_{\mathrm{ch}}(\cA)$ (the fully symmetric part) is the commutative chiral bar complex, and $F_1/F_2$ involves the Lie bracket (the antisymmetric part).  The associated graded $\mathrm{gr}^F \bar{B}_{\chirAss}$ is thus the Chevalley--Eilenberg complex for the Lie chiral algebra structure, which is $\bar{B}_{\chirCom}(\mathrm{Sym}^{\mathrm{ch}}(\cA))$ by the classical identification of the commutative bar complex with Chevalley--Eilenberg chains.

\emph{Step~3: Spectral sequence degeneration.}
For Koszul $\cA$, the spectral sequence $E_1^{p,q} = H^{p+q}(\mathrm{gr}^F_p \bar{B})$ degenerates at $E_1$.  To see this, note that the associated graded $\mathrm{gr}^F \bar{B}_{\chirAss}(\cA)$ is the Chevalley--Eilenberg complex (Step~2), and the Koszul property of $\cA$ as a chiral algebra implies that this commutative bar complex has cohomology concentrated on the diagonal $p + q = 0$ (i.e., $E_1^{p,q} = 0$ for $p + q \neq 0$).  Since $d_1\colon E_1^{p,q} \to E_1^{p+1,q}$ shifts the total degree $p+q$ by~$1$, it maps from the surviving diagonal to the zero region, hence $d_1 = 0$ and $E_2 = E_1$.  All higher differentials vanish for the same degree reason.  Therefore $H^*(\bar{B}_{\chirAss}(\cA)) \cong H^*(\mathrm{gr}^F \bar{B}_{\chirAss}(\cA))$, and the FG duality is recovered as the associated graded of our $\chirAss$ duality.

\emph{Step~4: Comparison with FG.}
The FG theorem states: for an $\Einf$-chiral algebra $\cA$, the chiral Koszul dual is a Lie chiral algebra $\cA^{!,\chirLie}$.  In the present framework: $\cA^{!,\chirAss} = \Omega(\bar{B}_{\chirAss}(\cA))$ is an $\Eone$-chiral algebra, and its associated graded under PBW is $\cA^{!,\chirLie} = \Omega(\bar{B}_{\chirCom}(\cA))$.  The FG duality is recovered.
\end{proof}

\section{Relationship to Ayala--Francis}
\index{Ayala--Francis}

\begin{tabular}{>{\raggedright\arraybackslash}p{0.38\textwidth}>{\raggedright\arraybackslash}p{0.52\textwidth}}
\textbf{Our Terminology} & \textbf{AF Terminology} \\
\hline
Non-abelian Poincar\'{e} duality & Nonabelian Poincar\'{e} duality (same) \\
FM compactification $\overline{C}_n(X)$ & Disk stratification of $\mathrm{Ran}(X)$ \\
Verdier duality on $\overline{C}_n(X)$ & Poincar\'{e}/Koszul duality for $\Eone$-algebras \\
Bar complex as factorization coalgebra & Factorization homology \\
\end{tabular}

\begin{remark}[Key comparison]\label{rem:ayala-francis-comparison}
The Ayala--Francis framework~\cite{AF15, AF20} establishes non-abelian
Poincar\'{e} duality as an equivalence between factorization homology and
factorization cohomology for framed $\mathsf{E}_n$-algebras on $n$-manifolds.
Our Theorem~A (Theorem~\ref{thm:bar-cobar-isomorphism-main}) is the
specialization to dimension~$n = 1$ (algebraic curves) with the additional
structure of:
\begin{enumerate}[label=(\roman*)]
\item \emph{Holomorphic enrichment.}
  The AF framework works with topological $\mathsf{E}_n$-algebras and
  locally constant cosheaves.  We work with holomorphic chiral algebras
  ($\mathcal{D}$-modules on curves), where the Fulton--MacPherson
  compactification provides the correct holomorphic compactification
  (not just the topological one).  The bar differential involves
  residues of meromorphic forms, not just topological data.
\item \emph{Explicit formulas.}
  The AF duality is an $\infty$-categorical equivalence; our realization
  via configuration space integrals gives explicit chain-level formulas
  for the bar differential (residues at collision divisors) and the
  cobar differential (inclusions of FM strata).
\item \emph{Higher genus.}
  For genus $g \geq 1$ curves, the configuration space
  $C_n(\Sigma_g)$ acquires additional topology ($H^1(\Sigma_g)$),
  and the bar complex inherits new generators from the handle
  cycles.  The AF framework predicts this (factorization homology
  on non-simply-connected manifolds), but the explicit identification
  with moduli space integrals and the Hodge bundle (Theorem~C)
  is new.
\end{enumerate}
\end{remark}


\section{Relationship to Feigin--Frenkel}
\index{Feigin--Frenkel}

\begin{tabular}{>{\raggedright\arraybackslash}p{0.38\textwidth}>{\raggedright\arraybackslash}p{0.52\textwidth}}
\textbf{Our Terminology} & \textbf{FF Terminology} \\
\hline
KM Koszul dual $\widehat{\fg}_{-k-2h^\vee}$ & FF involution $k \mapsto -k - 2h^\vee$ \\
Critical level $k = -h^\vee$ & Critical level (same) \\
Curvature $\kappa \propto (k + h^\vee)$ & Center of $\widehat{\fg}_{-h^\vee}$ \\
DS reduction $\mathcal{W}^k(\fg)$ & Drinfeld--Sokolov reduction (same) \\
$c + c' = 2d$ & Not formulated \\
Genus-$g$ obstruction $\mathrm{obs}_g$ & Not formulated \\
\end{tabular}

\begin{remark}[Key comparison]\label{rem:feigin-frenkel-comparison}
The Feigin--Frenkel involution $k \mapsto -k - 2h^\vee$ on affine
Kac--Moody levels~\cite{Feigin-Frenkel} is, in this monograph's framework,
a \emph{consequence} of chiral Koszul duality: the bar-cobar adjunction on
$\overline{C}_n(X)$ produces $\widehat{\fg}_{-k-2h^\vee}$ as the Koszul
dual of $\widehat{\fg}_k$ (Theorem~\ref{thm:universal-kac-moody-koszul}).
The geometric mechanism is Serre duality on the Fulton--MacPherson
compactification, which exchanges the chiral bracket
$\mu\colon j_*j^*(\cA \boxtimes \cA) \to \Delta_!\cA$
with the co-bracket on the Koszul dual.

At the critical level $k = -h^\vee$, the Feigin--Frenkel center
$\mathfrak{z}(\widehat{\fg}_{-h^\vee})$ is recovered as the degree-zero
homology of the (uncurved) bar complex: $\kappa = 0$ at $k = -h^\vee$,
so the bar complex is a genuine cochain complex, and
$H^0(\bar{B}(\widehat{\fg}_{-h^\vee})) \cong \mathfrak{z}(\widehat{\fg}_{-h^\vee})$
(Proposition~\ref{prop:km-generic-acyclicity}).

The extensions beyond FF are:
\begin{enumerate}[label=(\roman*)]
\item The complementarity formula $c(\widehat{\fg}_k) + c(\widehat{\fg}_{-k-2h^\vee}) = 2d$
  (Theorem~\ref{thm:central-charge-complementarity}), which shows the
  central charge sum is a root datum invariant, independent of the level.
\item The genus-$g$ obstruction formula $\mathrm{obs}_g = \kappa \cdot \lambda_g$
  (Theorem~\ref{thm:genus-universality}), which extends the FF involution
  to all genera via the Hodge bundle on $\overline{\mathcal{M}}_g$.
\item The Virasoro and $\mathcal{W}_N$ generalizations, where the
  Drinfeld--Sokolov parametrization introduces nonlinear level shifts
  (e.g., $c \mapsto 26 - c$ for Virasoro) that do not arise from a linear
  involution on levels.
\end{enumerate}
\end{remark}


\section{Relationship to Costello--Gwilliam}
\index{Costello--Gwilliam}

\begin{tabular}{>{\raggedright\arraybackslash}p{0.38\textwidth}>{\raggedright\arraybackslash}p{0.52\textwidth}}
\textbf{Our Terminology} & \textbf{CG Terminology} \\
\hline
Chiral algebra on $X$ & Holomorphic factorization algebra on $X$ \\
Bar complex $\B(\cA)$ & Factorization homology $\int_X \cA$ \\
Cobar construction $\Omega(C)$ & Factorization envelope \\
OPE residues & Observables in perturbative QFT \\
\end{tabular}

\begin{remark}[Key comparison]\label{rem:cg-comparison}
The Costello--Gwilliam factorization algebra framework~\cite{CG17}
relates to ours through the Beilinson--Drinfeld equivalence between
chiral algebras and factorization algebras on curves (BD, Chapter~3).
The precise comparison:
\begin{enumerate}[label=(\roman*)]
\item \emph{Factorization homology.}
  The CG factorization homology $\int_X \cA$ computes the same object
  as our chiral homology $H_*(\bar{B}^{\mathrm{ch}}(\cA))$.  Our explicit
  realization via configuration space integrals provides chain-level
  formulas that the CG framework treats abstractly.
\item \emph{Koszul duality.}
  The CG framework includes Koszul duality for factorization algebras
  (CG, Vol.~2, Chapter~5), specialized to holomorphic factorization
  algebras on curves.  Their Koszul duality, restricted to this setting,
  recovers our chiral Koszul duality.  Our contribution is the explicit
  identification of the Koszul dual for all standard families
  (Kac--Moody, Virasoro, $\mathcal{W}_N$, etc.) via the bar-cobar
  adjunction on Fulton--MacPherson spaces.
\item \emph{Quantum corrections.}
  The CG perturbative quantization framework produces $A_\infty$
  structures from Feynman diagram expansions (CG, Vol.~2, Chapter~3).
  Our curved $A_\infty$ structure on the bar complex is the
  chiral-algebraic counterpart: the curvature $m_0 = \kappa \cdot
  \mathbf{1}$ arises from the highest-order OPE pole (quartic for
  Virasoro, double for Heisenberg), not from a Feynman diagram sum.
  The genus universality theorem shows that these two perspectives
  produce the same genus-$g$ obstruction class.
\item \emph{BV-BRST.}
  The CG BV formalism produces the same BRST complex as our bar
  construction when applied to the chiral algebra of a 2d holomorphic
  field theory (Chapter~\ref{ch:bv-brst}).  The identification
  $\bar{B}(\cA) \simeq C^{\mathrm{BRST}}(\cA)$ connects the
  algebraic (bar complex) and physical (BRST cohomology) perspectives.
\end{enumerate}
\end{remark}

\section{Relationship to Loday--Vallette}

\begin{tabular}{>{\raggedright\arraybackslash}p{0.38\textwidth}>{\raggedright\arraybackslash}p{0.52\textwidth}}
\textbf{Our Terminology} & \textbf{LV Terminology} \\
\hline
Operad & Operad (same) \\
Koszul operad & Koszul operad (same) \\
Bar construction $\B$ & Bar construction $\B$ (same) \\
Cobar construction $\Omega$ & Cobar construction $\Omega$ (same) \\
Twisting morphism & Twisting morphism (same) \\
Maurer--Cartan equation & Maurer--Cartan equation (same) \\
$\Ass^! \cong \Ass$ & $\Ass^! \cong \Ass \otimes \mathrm{sgn}$ (Thm 7.1.1) \\
$\Com^! \cong \Lie$ & $\Com^! \cong \Lie\{1\}$ (Thm 7.2.1) \\
\end{tabular}
\begin{remark}[Convention for Koszul dual isomorphisms]
In Loday--Vallette \cite{LV12}, \S7.1--7.2, the Koszul dual includes the sign representation: $\Ass^! = \Ass \otimes \mathrm{sgn}$ and $\Com^! = \Lie\{1\}$ (with operadic suspension). In this manuscript, we absorb the sign twist into the bar desuspension: our bar construction uses the shifted generators $s^{-1}\mathcal{A}$ whose sign already accounts for the $\mathrm{sgn}$-twist. Concretely, our convention $\Ass^! \cong \Ass$ means that the dual cooperad is isomorphic to $\Ass$ as a \emph{graded} cooperad after this shift. All bar-cobar computations in the text use the desuspended convention consistently, so no sign errors propagate. See Appendix~\ref{app:signs} for the explicit translation between conventions.
\end{remark}

\begin{remark}[Chiral enhancement]
We lift the entire LV framework to the chiral setting, replacing vector spaces with D-modules and tensor products with chiral tensor products.
\end{remark}


\end{document}
